% !TeX program = xelatex
% Model specification, revised 2026-11-01. Compile three times with XeLaTeX.
\documentclass[11pt,a4paper]{article}
\usepackage[margin=2.35cm,headheight=15pt]{geometry}
\usepackage{amsmath,amssymb,bm,mathtools}
\usepackage{booktabs,tabularx,longtable,array}
\usepackage[no-config]{fontspec}
\usepackage{needspace,enumitem,fancyhdr}
\usepackage[unicode,hidelinks]{hyperref}
\defaultfontfeatures{Ligatures=Common}
\defaultfontfeatures[\rmfamily,\sffamily]{Ligatures=Common,Mapping=}
\setmainfont[AutoFakeSlant=0.2]{Noto Serif CJK SC}
\setsansfont{Noto Sans CJK SC}
\setmonofont{Noto Sans Mono CJK SC}
\XeTeXlinebreaklocale "zh"
\XeTeXlinebreakskip=0pt plus 1pt minus 0.1pt
\setlength{\parindent}{2em}
\setlength{\parskip}{0.35em}
\setlength{\emergencystretch}{3em}
\setlength{\LTpre}{0.6em}
\setlength{\LTpost}{0.6em}
\renewcommand{\arraystretch}{1.22}
\setlist{itemsep=0.2em,topsep=0.4em,leftmargin=2em}
\allowdisplaybreaks[2]
\setcounter{tocdepth}{1}
\setcounter{secnumdepth}{3}
\numberwithin{equation}{section}
\renewcommand{\contentsname}{目录}
\renewcommand{\refname}{参考文献}
\renewcommand{\tablename}{表}
\pagestyle{fancy}
\fancyhf{}
\fancyhead[L]{\small Model 1031：经验、概括与规则}
\fancyhead[R]{\small 2026-11-01}
\fancyfoot[C]{\thepage}
\newcolumntype{L}[1]{>{\raggedright\arraybackslash}p{#1}}
\newcolumntype{Y}{>{\raggedright\arraybackslash}X}
\newcommand{\ind}{\mathbb{I}}
\newcommand{\norm}[1]{\left\lVert#1\right\rVert}
\newcommand{\clip}{\operatorname{clip}}
\newcommand{\logit}{\operatorname{logit}}
\newcommand{\sigmoid}{\operatorname{sigmoid}}
\newcommand{\A}{\mathcal{A}}
\newcommand{\Hh}{\mathcal{H}}
\newcommand{\Rr}{\mathcal{R}}
\newcommand{\Dd}{\mathsf{D}}
\newcommand{\Ll}{\mathsf{L}}
\newcommand{\Ee}{\mathsf{E}}
\newcommand{\Ii}{\mathcal{I}}
\newcommand{\Tt}{\mathcal{T}}
\DeclareMathOperator*{\argmin}{arg\,min}
\DeclareMathOperator*{\argmax}{arg\,max}
\hypersetup{pdftitle={Model 1031：经验、概括与规则的学习模型},pdfauthor={},bookmarksdepth=3}
\title{Model 1031\\[0.4em]\Large 具体经验、局部概括与规则的共同学习模型}
\author{}\date{2026年11月1日}
\begin{document}
\maketitle
\thispagestyle{empty}

\noindent\textbf{本轮定稿：十参数共同学习模型。}

Model 1031把具体样例、局部概括和目录规则放在同一个有限工作区中。
学习者先保存完成的经历，再尝试忽略部分特征，形成适用范围更广的概括；
这些概括又为目录规则提供进入机会。每项在场内容都预测类别，实际反馈统一评价这些预测，
并改变它们以后影响选择的程度。
旧样例、概括和规则可以同时保留，形成过程允许反复保存新样例和替换已有内容。

该设计把表示形式与学习机制分开。具体样例保存局部经验，规则表达可推广的类别关系；
两者共享容量、内容信念、反馈更新和选择接口。它们在每题的总影响由当前内容及其权重决定。
判断精度随各项内容解释反馈的表现而调整；内容信念与执行选择分别定义。
旧信念经规定的保持步骤后，与本题反馈提供的证据共同决定下一题的信任程度。
因此，经验与规则之间的关系同时包含内容衍生和当前选择中的概率混合。

规则与具体记忆共同参与分类已有RULEX等过程模型的理论基础\cite{nosofsky1994}。
有限采样也被用于解释类别学习中的近似推断过程\cite{sanborn2010}。
本模型把有限工作区、逐步概括和共同反馈评价规定为具体的状态转移。
这些形成操作给出可检验的行为预测，模拟恢复检查推断方法在已知模型下能识别哪些状态。

本文给出三种反馈条件下的统一实现。条件3同时保存内容、类别和同科配对之间的关联，
使半对反馈能够直接约束联合信念。本轮先用开发样本比较八种版本，再以保留样本检查初选方案。七参数初选方案未通过检验，
随后用已冻结的训练参数比较九参数方案与完整版本，并据此确定本轮采用版本。
该修订比较属于开发分析；第\ref{sec:evidence}节记录原检验、修订决定及恢复限制。

{\small\setlength{\parskip}{0pt}\tableofcontents}
\clearpage
\section{任务、状态与逐题顺序}
\label{sec:state}

本版的任务语义对应exp123的Task2。刺激有四个特征，每个特征已归一到$[0,1]$。
用$\bm x_t$表示第$t$题的物理刺激，
$a_t\in\{1,\ldots,C\}$表示选择，$r_t$表示反馈。条件1有两个类别，条件2、3有四个类别。
条件1、2的反馈为正确或错误；条件3的反馈为种正确、同科但种错误、科错误，数值分别为
$1,\tfrac12,0$。选择编号与按键的对应是已知记录信息。正确类别和真实同科配对用于实验
环境产生反馈，学习者接收刺激、自己的选择及反馈。
本批条件3的环境编码将类别1、2归为一科，将类别3、4归为另一科。
模型从正式任务开始，按记录顺序连接同一被试的各场次。场次间工作区与信念延续，
休息时长未作为状态转移的输入。

知觉模块先产生内部刺激$\bm u_t$。误差分布沿用共享模块的被试分组设置，参数来自独立
知觉任务的既有汇总。正态误差使用各维的固定均值$\mu_d^{\mathrm P}$和标准差
$\sigma_d^{\mathrm P}$；均匀误差使用知觉阈值给出的半宽$\tau_d^{\mathrm P}$：
\begin{equation}
\begin{aligned}
u_{td}&=\clip_{[0,1]}(x_{td}+\epsilon_{td}),\\
\epsilon_{td}&\sim\begin{cases}
\mathcal N\!\left(\mu_d^{\mathrm P},(\sigma_d^{\mathrm P})^2\right),&\text{正态误差组},\\
\operatorname{Uniform}(-\tau_d^{\mathrm P},\tau_d^{\mathrm P}),&\text{阈值均匀误差组}.
\end{cases}
\end{aligned}
\label{eq:perception}
\end{equation}
各维及各题的误差独立抽取；低于0的值设为0，高于1的值设为1。
本轮S128、S132、S229、S331采用正态误差，其余14人采用阈值均匀误差。
各人的实际分布类型及按特征顺序排列的固定数值保存在研究目录的
\path{perception_configuration_audit.json}。
同一知觉实现用于行为推断和自主生成。
这些知觉参数在本轮固定，其独立任务估计误差尚未传播到类别学习参数及状态分布中。
知觉误差与内容搜索、选择抽样分别使用随机数流。它们使固定参数也能产生多种学习路径。

特征按输入的feature1--feature4顺序保存，各维的身体部位名称由个人任务设置给出。

工作区最多容纳$K$项内容。当前活跃集合记为$\mathcal H_t$，内容$h$的信念为$\pi_{th}$，
满足$\pi_{th}\geq0$及$\sum_{h\in\mathcal H_t}\pi_{th}=1$。这里的信念表示当前保留内容之间
的相对可信度。工作区保存以下两种形式，三种心理内容由其具体状态区分。

\begin{center}
\begin{tabularx}{\textwidth}{p{2.0cm}Y Y}
\toprule
内容 & 保存的信息 & 对新刺激的预测\\
\midrule
具体样例 & 四维中心$\bm v_h$、全部四个维度、类别信念 & 在全部维度上相似时支持记住的类别\\
局部概括 & 中心$\bm v_h$、保留维度集合$M_h$、类别信念 & 只在保留维度上比较相似程度\\
目录规则 & 固定目录中的规则编号 & 通过类别区域和边界距离预测类别\\
\bottomrule
\end{tabularx}
\end{center}

具体样例和局部概括共同采用经验模板的存储形式。$M_h$含四个维度时是具体样例；
$1\leq |M_h|<4$时是局部概括。每项内容另有路径内唯一编号、来源编号及进入试次，
用于追踪形成与替换。存储位置可重复使用，不同模拟路径上的位置和编号分别解释。

一题开始时，模型先产生$\bm u_t$，再根据已经完成的经历决定是否加入新内容。
随后形成选择分布并产生$a_t$。接收$r_t$后，模型用更新前的预测计算反馈似然，
更新内容信念和模板类别信念，并保存本题经历。该经历从下一题开始可参与内容形成。
首次有效反馈后的下一题，模型确定保存第一项具体样例；第一题使用均匀选择。
为统一初始题的记号，工作区先保存一个“尚未形成内容”的占位状态，其信念为1、
各类别预测为$1/C$，条件3的三种配对等权。首个样例进入时替换该状态。
初始占位与具体样例、局部概括、目录规则分别统计。

\section{共同的类别预测接口}
\label{sec:emission}

每项内容都把内部刺激变为一个类别概率向量。这个共同接口使具体样例、局部概括和规则
能够接受同一次反馈的评价。以下公式中的$h$均指当前路径中确实在场的一项内容。
单题公式省略试次下标$t$，上标$+$表示完成当前反馈更新之后的量。

\subsection{经验模板的相似度与类别信念}
模板只比较仍被保留的维度。其相似度，也称亲和度（affinity），定义为
\begin{equation}
A_h(\bm u)=\exp\!\left[-\frac{\sum_{d\in M_h}(u_d-v_{hd})^2}{2\sigma^2}\right].
\label{eq:affinity}
\end{equation}
$\sigma>0$规定相似范围。$A_h$在0和1之间，完全匹配时为1。删除一个保留维度后，
该维度的差异不再降低相似度，因此模板能够覆盖更多刺激。中心$\bm v_h$保持为形成时的
内部刺激；本版通过删除维度形成概括。

模板还保存所关联类别的不确定性。条件1、2用$\theta_h(z)$表示模板关联类别$z$的概率，
各类别概率之和为1。判断精度$\beta_h$控制匹配后对该类别的偏向程度。定义
\begin{equation}
s_C(\beta)=\frac{1-e^{-\beta}}{1+(C-1)e^{-\beta}},\qquad
b_h(\bm u)=A_h(\bm u)s_C(\beta_h).
\label{eq:gain}
\end{equation}
相似度和判断精度相乘形成有效增益$b_h$。假定模板关联类别为$z$时，其类别预测为
\begin{equation}
E_h(c\mid z,\bm u)=\frac{1-b_h(\bm u)}{C}+b_h(\bm u)\ind(c=z).
\label{eq:template-component}
\end{equation}
$\ind(\cdot)$为指示函数，条件成立时取1。对未知的关联类别求和，可得模板实际使用的预测：
\begin{equation}
p_h(c\mid\bm u)=\sum_z\theta_h(z)E_h(c\mid z,\bm u)
=\frac{1-b_h(\bm u)}{C}+b_h(\bm u)\theta_h(c).
\label{eq:template-prediction}
\end{equation}
这一形式先把关联类别得分设为$\beta_h$、其余类别得分设为0，通过softmax得到完全匹配时的
类别预测；再以相似度$A_h$为比例，与均匀预测混合。对未知关联类别求平均后，得到上式。
刺激相距很远时，$b_h$接近零，预测接近均匀分布。刺激相似且判断精度高时，预测接近
模板的类别信念。$\sigma$控制适用范围，$\beta_h$控制匹配时的判别强度，$\theta_h$
表示关联类别的不确定性。这三个量的心理职责不同，其估计仍可能相互补偿。

\subsection{目录规则的边界预测}
规则目录规定每条规则$h$对应的类别区域$R_{hc}\subseteq[0,1]^4$。令
$d_{hc}(\bm u)=\inf_{\bm z\in R_{hc}}\|\bm u-\bm z\|_2$为刺激到该类别区域的欧氏距离。
规则的预测是距离的softmax归一化：
\begin{equation}
p_h(c\mid\bm u)=\frac{\exp[-\beta_h d_{hc}(\bm u)]}
{\sum_{j=1}^C\exp[-\beta_h d_{hj}(\bm u)]}.
\label{eq:rule-prediction}
\end{equation}
刺激处于某类别区域内时，到该区域的距离为零；距离较远的类别得到较低概率。归一化
使各类别概率之和为1。实现直接调用共享边界几何及其精确活跃集投影算法，数值容差为$10^{-9}$。

目录把每个类别区域存成有限个凸多面体的并集：
\begin{equation}
R_{hc}=\bigcup_{\ell}\{\bm v\in[0,1]^4:A_{hc\ell}\bm v\leq\bm b_{hc\ell}\}.
\label{eq:catalogue-regions}
\end{equation}
矩阵$A_{hc\ell}$的每一行和向量$\bm b_{hc\ell}$的对应元素共同规定一个线性约束，
同时满足这些约束的点构成一个凸区域。取并集可以表达“接近”区域之外的两侧等非凸类别。
区域系数和类别标签由固定目录给定，距离计算对这些区域进行投影。

目录保留29条二分类规则和152条四分类规则，包含单维判断、维度间比较、和与差的关系、
局部分支等预设表达资源。目录的关系运算及固定类别标签属于模型假设。
本版研究这些表达如何进入工作区、如何获得信任以及何时影响选择。目录覆盖及标签对称性
另行记录。现有目录没有包含每个几何结构的全部类别标签排列，任意反应重编号会改变可用
规则预测及其先验。当前个体参数与内容推断以该固定目录和\texttt{choice}编码为条件；
按键标签的完整学习过程仍需单独检验。
对本轮18名被试的9792个记录刺激，二分类目录中的规则0、四分类目录中的规则42均能
正确分类各自条件的全部记录，编号从0开始。此检查只读取物理刺激与实验类别，界定
目录对记录日程的覆盖。学习者发现、相信及执行该规则的概率由后续学习过程决定。

\subsection{条件3的配对条件预测}
条件3中，学习者知道四个选项组成两科、每科两种，具体配对关系需要学习。
可能配对记为$g\in\mathcal G=\{12|34,13|24,14|23\}$。
每项内容保存条件配对信念$q_h(g)$；每个模板进一步保存
$\theta_{hg}(z)$，即给定配对$g$时模板关联类别$z$的概率。
模板使用
\begin{equation}
p_{hg}(c\mid\bm u)=\frac{1-b_h(\bm u)}4+b_h(\bm u)\theta_{hg}(c).
\label{eq:conditional-template}
\end{equation}
规则的类别几何与$g$无关，因而$p_{hg}=p_h$。每项内容的类别边缘预测为
$p_h(c\mid\bm u)=\sum_gq_h(g)p_{hg}(c\mid\bm u)$。
内容与配对的联合信念写为$J_{hg}=\pi_hq_h(g)$，全体$h,g$上的概率总和为1。
各$q_h$描述同一个任务配对在不同内容成立时的条件信念；任务配对的总信念为
$\sum_h\pi_hq_h(g)$。内容与配对之间的相关性由$J$保留。
数学上可保存$J$并由它计算$\pi$与$q$；后两者是边缘与条件分布。实现缓存这些派生量，
以复用共享接口，并在更新和状态恢复后保持它们一致。
条件3的每个模板保存$3\times4$个条件类别概率，每项内容保存3个配对概率，均满足相应的
归一化约束。这些量随经历更新。三条件采用相同的自由参数定义与维数，本轮分别估计
每个被试的参数。

配对与类别的关联对保留半对信息很重要。例如，第一只动物选择A得到半对，模型保存三种
可能：AB同科且动物属于B，AC同科且动物属于C，AD同科且动物属于D。
初始配对边缘仍可各为$1/3$，每种配对内的类别却已经受到约束。
在理想记忆和精度极限下，再次对同一动物选择A会预测半对，其概率趋于1。
此时，两个独立的均匀边缘会给出$1/3$。条件类别分布保留了这项反馈提供的关联信息。
如果下一次匹配这个模板时，选择B得到完全正确，AB配对的相对信念就会提高。
具体样例因此可以在目录规则形成之前，利用前后两次反馈支持配对学习。
作为有限精度的例子，假设当前只有这项模板、三种配对等权，且匹配增益$b_h=0.9$。
选择B并得到完全正确时，AB配对的反馈概率为$0.925$，另外两种配对均为$0.025$。
于是AB配对更新为$0.925/(0.925+0.025+0.025)\approx0.949$。
这一变化来自两次经历共同提供的类别与配对约束。

\section{从经历形成内容}
\label{sec:formation}

新内容的形成使用已经完成的经历。搜索负责提出待检验的内容，后续真实反馈负责评价它们。
保存具体样例时，使用紧邻的上一题内部刺激、选择及反馈；已经保留的模板继续作为概括和
规则提议的来源。
本节规定内容形成和容量规则；控制搜索机会的参数见
第\ref{sec:selected}节。

\subsection{保存样例与逐步概括}
首个具体样例在第一次有效反馈后的下一题开始时进入工作区，保存上一题的内部刺激，
保留全部四个维度。
条件1、2的反馈直接界定这次经历允许的类别：正确时关联已选类别，错误时在其余类别中
均匀分配信念。四分类中选择A后得到错误，新样例记录
$(0,1/3,1/3,1/3)$。该初始化使用来源经历的信息。

条件3对每种配对分别确定允许类别。正确反馈保留所选类别，半对保留该配对下的同科另一个
类别，全错保留另一科的两个类别。模板在各配对的允许集合内均匀初始化$\theta_{hg}$，
并从现有工作区继承配对边缘$\sum_j\pi_jq_j(g)$。初始工作区的配对信念为均匀分布。
各配对下允许集合大小相同，故该初始化保留继承的配对权重。

保存更多样例和形成局部概括构成两种局部操作。存在可删维模板时，两种操作各分配一半机会；
形成概括时，先按内容信念从可删维模板中抽取父模板，再均匀删除其一个保留维度，至少保留
一个维度。各形成操作在可选父项总信念为零时，均采用均匀抽取。新概括继承中心及完整类别信念；
条件3同时继承父模板的条件配对信念。
父模板和新概括可以同时保留，并在之后分别接受反馈。

例如，一个模板原本比较头、脖子、腿和尾巴。删除尾巴维度后，尾巴不同的动物也可能与它
高度匹配。后续反馈一致时，该概括的内容信念提高；反馈冲突时，相对信念下降。
概括因此可以形成、保留或被反例削弱。

\subsection{从概括提出目录规则}
工作区出现局部概括后，以$g_{\mathrm R}$的概率尝试目录规则。先按内容信念选择一个父概括，
再从尚未在场的规则中抽取候选。若这次规则尝试没有可用规则，则继续局部操作。
候选规则的概率由其与父概括类别倾向的一致程度决定。

为评价这种一致程度，保留维度取父模板中心的值，被忽略的维度分别取0和1，组合成检查点集
$\mathcal P_h$。至多三个维度被忽略，因此检查点至多八个。
令$\bar\theta_h=\theta_h$用于条件1、2，
$\bar\theta_h(z)=\sum_gq_h(g)\theta_{hg}(z)$用于条件3。
令$\mathcal R_t$为当前未在工作区中的目录规则集合。候选规则$k\in\mathcal R_t$的
提议分数和概率为
\begin{equation}
\begin{split}
Q(k\mid h)&=\exp\!\left\{\frac1{|\mathcal P_h|}
\sum_{\bm v\in\mathcal P_h}\log\left[\sum_c\bar\theta_h(c)p_k(c\mid\bm v)\right]\right\},\\
P(k\mid h,\mathrm{proposal})&=\frac{Q(k\mid h)}{\sum_{k'\in\mathcal R_t}Q(k'\mid h)}.
\end{split}
\label{eq:proposal}
\end{equation}
这里用父概括的当前判断精度计算候选规则预测。括号内是规则预测与父模板类别信念的重合程度，
对检查点取几何平均，使在这些位置持续一致的规则获得较高分数，并使不同检查点数采用相同的分数尺度。
检查点为规则提议提供有限近似，后续真实反馈继续检验其适用范围。
这些检查点只用于形成提议。类别信念的反馈更新使用真实完成的试次。
条件3的新规则继承父概括的$q_h$，类别预测由自己的边界给出。

\subsection{容量、替换和入场信念}
每题最多提出一个新候选。候选与在场内容重复时，本题沿用原工作区，下一题仍可搜索。
形成与替换均使用上一题反馈之后的内容信念；加入后归一化为当前题选择前信念。
规则重复按规则编号判断；模板重复按保留维度、这些维度上的中心及完整类别信念判断。
重复时保留原项及其配对信念。这些去重键规定了模型采用的内容粒度。

工作区未满且已有$n$项内容时，新内容获得$\mu=1/(n+1)$的信念。
工作区已满时，删除信念最低的一项，并列者随机抽取。令其原信念为$\pi_e$，
$y_{t-1}=\ind(r_{t-1}=1)$表示上一题完全正确，则入场信念为
\begin{equation}
\mu=y_{t-1}\pi_e+(1-y_{t-1})/K.
\label{eq:admission-mass}
\end{equation}
正确后的新内容继承被替换项的质量；未完全正确后的新内容获得均匀份额$1/K$。
其余内容按原相对比例分享$1-\mu$。首个样例的信念为1；$K=1$时，新内容完整替换原内容。
完全正确时，幸存内容的绝对信念因此保持原值。未完全正确时，由$\pi_e\leq1/K$可知，
新候选获得的份额至少与被替换项相同，使新解释有机会接受后续检验。
条件3的幸存内容保留各自$q_h$和$\theta_{hg}$，新内容采用相应来源的继承状态。

\section{共同的反馈学习}
\label{sec:feedback}

反馈似然回答一个共同问题：如果当前内容的类别预测成立，实际看到的反馈有多大概率发生？
定义$T_g(r\mid a,c)$为反馈事件的指示函数：若正确类别为$c$、选择为$a$且配对为$g$，
实验会给出$r$，则取1。条件1、2省略$g$。给定内容和配对，反馈似然为
\begin{equation}
L_{hg}=\sum_{c=1}^C T_g(r_t\mid a_t,c)p_{hg}(c\mid\bm u_t).
\label{eq:feedback-likelihood}
\end{equation}
每项在场内容都计算这项似然。当前选择确定了实际反馈事件，反馈随后用于评价整个工作区
中的内容。
该似然等于反馈允许的类别概率之和。条件1、2的正确似然为$p_h(a_t)$，错误似然为
$\sum_{c\ne a_t}p_h(c)$。条件3令$s_g(a)$为配对$g$下与$a$同科的另一个选项，得到
\begin{equation}
L_{hg}=\begin{cases}
p_{hg}(a_t),&r_t=1,\\
p_{hg}(s_g(a_t)),&r_t=\tfrac12,\\
\displaystyle\sum_{c\notin\{a_t,s_g(a_t)\}}p_{hg}(c),&r_t=0.
\end{cases}
\label{eq:ternary-feedback}
\end{equation}
三种事件互斥且概率总和为1。半对反馈通过其所限定的类别集合进入似然。
在配对均匀且类别预测与配对无关时，半对似然为$(1-p_h(a_t))/3$，全错似然是其两倍。
该条件下两者提供相同的相对内容证据。配对知识或条件类别信念形成后，这一等价性可以消失。

\subsection{信念保持与联合反馈更新}
用$\widetilde\pi_h$和$\widetilde q_h(g)$表示保持步骤后的有效旧信念。
令$y_t=\ind(r_t=1)$，内容保持指数取$\delta_t=\delta+(1-\delta)y_t$。
完全正确时保留原内容信念；未完全正确时，较小的指数减弱原有优势。
配对信念每题使用同一个参数$\delta$，保持作用于内容边缘及各内容的条件配对分布。
对保持率$\delta\in[0,1]$，模型通过幂归一化实现旧证据保持：
\begin{equation}
\widetilde\pi_h=\frac{\pi_h^{\delta_t}}{\sum_j\pi_j^{\delta_t}},\qquad
\widetilde q_h(g)=\frac{q_h(g)^\delta}{\sum_kq_h(k)^\delta}.
\label{eq:retention}
\end{equation}
指数小于1时，原本集中的分布变得平缓，旧信念的优势减弱。
内容边缘与每项内容内部的配对分布分别归一化，使配对的不确定性通过实际反馈似然影响
内容信任。对零质量条件分布，数值实现使用均匀回退；该分支以其内容质量参与预测。
保持指数为零时，相应有效旧分布定义为当前活跃支持上的均匀分布。

保持之后的联合分布为$\widetilde J_{hg}=\widetilde\pi_h\widetilde q_h(g)$。
条件3将它与实际反馈的似然相乘并归一化：
\begin{equation}
J^+_{hg}=\frac{\widetilde\pi_h\widetilde q_h(g)L_{hg}}
{\sum_{j\in\mathcal H_t}\sum_{k\in\mathcal G}\widetilde\pi_j\widetilde q_j(k)L_{jk}},\qquad
\pi_h^+=\sum_gJ^+_{hg},\qquad q_h^+(g)=\frac{J^+_{hg}}{\pi_h^+}.
\label{eq:joint-update}
\end{equation}
分母对全部活跃内容和配对求和，使联合概率归一化。条件1、2对应只有一个配对状态的特例，
因此退化为$\pi_h^+\propto\widetilde\pi_hL_h$。

\subsection{模板类别的条件更新}
模板还需要学习它关联哪个类别。对一个可能的关联类别$z$，用
\begin{equation}
D_{hgz}=\sum_cT_g(r_t\mid a_t,c)
\left[\frac{1-b_h(\bm u_t)}C+b_h(\bm u_t)\ind(c=z)\right]
\label{eq:label-likelihood}
\end{equation}
计算反馈概率。$b_h$使用当前反馈之前的判断精度。于是
$L_{hg}=\sum_z\theta_{hg}(z)D_{hgz}$，而条件类别信念更新为
\begin{equation}
\theta^+_{hg}(z)=\frac{\theta_{hg}(z)D_{hgz}}{L_{hg}}.
\label{eq:label-update}
\end{equation}
条件1、2省略$g$。模板的类别信念按当前反馈条件更新。
新内容的来源反馈用于初始化，从加入后的真实试次开始接受可信度评价。

条件1的一次反馈唯一确定样例类别，因此新模板的$\theta_h$只有一个分量为1，其余为0。
按式\eqref{eq:label-update}更新后，该确定类别保持不变；后续反馈改变模板的内容信念。
条件2的错误反馈、条件3在给定配对下的全错反馈允许多个类别，条件类别更新继续调整这些
类别的相对信念。条件3的半对反馈在每种配对内限定一个类别，类别边缘的不确定性
同时随配对信念更新。

对于模板内容，这些更新对应保持步骤之后的同一次联合条件化。把式\eqref{eq:joint-update}与
式\eqref{eq:label-update}相乘，可见
\begin{equation}
J^+_{hg}\theta^+_{hg}(z)\ \propto\
\widetilde\pi_h\widetilde q_h(g)\theta_{hg}(z)D_{hgz}.
\label{eq:one-event}
\end{equation}
右侧只含一次反馈事件概率$D_{hgz}$。
内容可信度、配对关系和关联类别分别是同一个联合分布的边缘与条件信息。
规则内容直接使用式\eqref{eq:joint-update}更新，其质量与模板部分通过同一个分母共同归一化。

\section{最终保留的控制机制与参数}
\label{sec:selected}
本轮采用十参数版本，比较标识为\texttt{full}。
其参数向量为
\begin{equation}
\boldsymbol\psi=(K,\sigma,\beta_0,\alpha,p_0,\eta,\rho,g_{\mathrm R},\delta,\kappa).
\label{eq:parameters}
\end{equation}
容量$K$决定同时保存多少项内容，相似宽度$\sigma$决定模板覆盖范围，初始判断精度$\beta_0$和学习率$\alpha$
决定类别预测的集中程度及其随反馈的变化。保持率$\delta$控制旧内容与配对信念的保持。三个搜索参数$p_0,\eta,\rho$控制何时尝试新内容，
$g_{\mathrm R}$控制已有概括为目录规则提供的尝试机会。
它用于内容形成；作答阶段的各类内容权重由执行概率汇总。
$\kappa$控制继续执行上一题内容的额外倾向。

\begin{center}
\begin{tabularx}{\textwidth}{p{1.2cm}p{2.3cm}Y}
\toprule
参数 & 数学范围 & 本轮搜索取值\\
\midrule
$K$ & 正整数 & 1、2、3、4、6、8、10\\
$\sigma$ & $>0$ & 0.08、0.15、0.25、0.30、0.40、0.60、0.65、1.00\\
$\beta_0$ & $[0.1,25]$ & 0.5、1.5、2、3、5、6、10、12、25\\
$\alpha$ & $\geq0$ & 0、0.01、0.03、0.04、0.10、0.30\\
$p_0$ & $(0,1)$ & 0.01、0.02、0.03、0.08、0.15、0.25、0.30、0.50\\
$\eta$ & $\geq0$ & 0、0.10、0.25、0.50、0.70、1.00、1.50、3.00\\
$\rho$ & $[0,1]$ & 0、0.30、0.40、0.70、0.80、0.90、0.95、0.98\\
$g_{\mathrm R}$ & $[0,1]$ & 0.03、0.08、0.10、0.20、0.25、0.30、0.50、0.70、0.80\\
$\delta$ & $[0,1]$ & 0.40、0.65、0.70、0.80、0.90、0.97、1.00\\
$\kappa$ & $\geq0$ & 0、0.5、1、1.5、3、6\\
\bottomrule
\end{tabularx}
\end{center}
数学范围规定模型允许的参数，表中离散取值合并了新搜索网格与16个既有预设点的坐标。
本轮采用有限候选组合，参数取值的笛卡尔积并未全部搜索。
知觉参数由既有独立任务汇总给定，局部操作各一半的比例及目录几何作为固定模型假设。

\subsection{失败累积控制搜索机会}
令$F_t$为第$t$题开始时的累积失败量，$F_1=0$，并令$y_t=\ind(r_t=1)$。
完成反馈后，该量按
\begin{equation}
F_{t+1}=\begin{cases}\rho F_t,&y_t=1,\\ F_t+1,&y_t=0\end{cases}
\label{eq:search-pressure}
\end{equation}
变化。因此连续未完全成功会逐步增加累积失败量，完全成功后该量按比例降低。
下一题出现一次新内容尝试的概率为
\begin{equation}
p^{\mathrm{search}}_{t+1}=\frac{1}{1+\exp\{-\log[p_0/(1-p_0)]-\eta F_{t+1}\}}.
\label{eq:search-probability}
\end{equation}
$p_0$给出无累积失败时的基准机会，$\eta$决定失败对搜索机会的影响，$\rho$决定成功后
此前累积失败量保留多少。条件3的半对与全错都使$F$增加1，完整三值反馈同时进入
式\eqref{eq:ternary-feedback}的证据更新。第一次有效经历确定形成首个样例。
之后的搜索尝试再按第\ref{sec:formation}节生成样例、概括或规则，并接受去重与容量检查。
当$\eta=0$时，$\rho$对行为预测失去作用；本轮搜索将该边界统一记录为$\rho=0$，
并在参数恢复中把它标为无效坐标。

\subsection{精度学习与证据保持}
各项内容的判断精度从共同初值$\beta_0$出发。每次反馈后，先用更新前的预测计算实际反馈的
对数似然，再沿该似然对精度的梯度更新：
\begin{equation}
\beta_h^+=\operatorname{clip}_{[0.1,25]}
\left[\beta_h+10\alpha\frac{\partial\log L_h}{\partial\beta_h}\right].
\label{eq:selected-precision}
\end{equation}
$\alpha$为精度学习率，固定倍数10沿用本版精度更新约定。
条件3中$L_h=\sum_g\widetilde q_h(g)L_{hg}$，求导时固定配对和类别信念。
新内容继承加入之前的信念加权平均精度。内容、配对和模板类别的当前反馈更新均使用
精度改变之前的预测，保持同一反馈事件的计算时点一致。

反馈调节的内容保持指数为$\delta_t=\delta+(1-\delta)y_t$。
完全正确时旧内容信念完整进入更新，未完全正确时按$\delta$作幂归一化。
条件配对信念每题按同一个参数$\delta$保持，模板类别信念继续按实际反馈条件更新。
幂保持后的有效旧信念与当前反馈似然相乘并归一化，见式\eqref{eq:joint-update}。
这一认知规则对旧证据的使用依赖反馈类型，其整体解释由明确的状态转移模型规定。

\subsection{信念与执行维持共同决定选择}
内容信念表示对在场内容的相对信任。执行维持参数$\kappa\geq0$表示继续使用上一题内容的
额外倾向。令$H_{t-1}$为上一题实际执行内容的唯一身份，则本题的执行分布为
\begin{equation}
w_{th}=\frac{\pi_{th}\exp\{\kappa\ind(h=H_{t-1})\}}
{\sum_{j\in\mathcal H_t}\pi_{tj}\exp\{\kappa\ind(j=H_{t-1})\}},
\qquad H_t\sim\operatorname{Categorical}(w_t).
\label{eq:execution-policy}
\end{equation}
例如，两项内容的信念均为$1/2$，上一题使用第一项；$\kappa=\log3$时，继续使用它的概率为$3/4$。
执行维持因此可以让相同信念产生不同的连续使用模式。学习仍由所有在场内容共同接受反馈。
上一题内容被替换后，该身份结束，各幸存或新增内容按当前信念重新取得执行机会。

当$\kappa=0$时，$w_{th}=\pi_{th}$，给定实际选择与反馈，旧执行身份对未来状态的影响消失。
该边界可在行为似然中逐题解析求和执行身份。$\kappa>0$时，实际执行身份影响下一题，
行为推断保留其路径历史，自主生成与推断使用同一执行转移。

\subsection{足以推进下一题的状态}
工作区结构$\mathcal H_t$保存各模板的中心、保留维度及各目录规则的编号。
$J_t$保存内容与配对的联合信念，$\Theta_t$保存模板的条件类别信念，
$\boldsymbol\beta_t$保存各内容的判断精度，$F_t$保存累积失败量。
在内容形成之后、抽取执行内容之前，用
\begin{equation}
S_t=(\mathcal H_t,J_t,\Theta_t,\boldsymbol\beta_t,F_t,H_{t-1})
\label{eq:minimal-state}
\end{equation}
规定本题执行分布。若上一题内容已不在场，$H_{t-1}$记为空身份；第一题也使用空身份。
内部刺激$\bm u_t$经共同预测接口变为各内容的类别预测。
抽到$H_t$后，旧身份的作用已传递给本题执行身份，接下来采用以下状态推进路径：
\begin{equation}
Z_t=(\mathcal H_t,J_t,\Theta_t,\boldsymbol\beta_t,F_t,\bm u_t,H_t).
\label{eq:path-state}
\end{equation}
给定$Z_t$，回答从已选内容的类别分布抽取。完成本题的内部刺激、回答和反馈共同用于
产生下一题的反馈学习、知觉与内容形成，再抽取新的执行内容。

条件1、2把配对$g$的支持取为单点，此时$J_t=\pi_t$；条件3由$J_t$计算内容边缘$\pi_t$
和条件配对$q_t$。这些边缘与条件量是联合状态的派生信息。
固定参数与$Z_t$足以规定后续状态转移；实现另外保存来源编号、日志与概率缓存，用于追踪
完整路径和访问数组。

\section{选择、贝叶斯推断与自主生成}
\label{sec:inference}

\subsection{选择分布及两类权重}
本题采用内容$h$的概率$w_{th}$由内容信念与上一题执行身份共同决定，见式\eqref{eq:execution-policy}。
$S_t$表示本题内容形成之后、执行内容抽样之前的状态。
给定该状态和当前内部刺激，选择概率为
\begin{equation}
P(a_t=c\mid S_t,\bm u_t)=\sum_{h\in\mathcal H_t}w_{th}
\sum_{g\in\mathcal G}q_{th}(g)p_{thg}(c\mid\bm u_t).
\label{eq:choice}
\end{equation}
条件1、2省略配对求和。把$w_{th}$对所有具体样例相加，得到本题采用具体样例的总概率；
对局部概括或目录规则分别求和也得到相应总概率。三类内容的影响随形成、替换和反馈而变化。
它们通过来源关系和共同信念相互联系。首个样例形成后，三类总权重之和为1；
初始题的概率质量由上述占位状态承担。

例如，某题三类内容的总执行权重依次为0.2、0.3和0.5，组内加权预测选择A的概率分别为
0.8、0.6和0.4，则总选择概率为$0.2\times0.8+0.3\times0.6+0.5\times0.4=0.54$。
这是加权计算的示例；每题的组权重由当前在场内容的执行权重相加得到，并随学习过程变化。

学习者的内容信念$\pi_{th}$与研究者的粒子权重$W_t^{(i)}$有不同含义。
$\pi$描述一条认知路径内对内容的信任；$W$描述观察到被试行为后，研究者对不同潜在路径的
相对支持。条件配对信念$q_h$属于前者。粒子滤波按观察到的选择更新研究者的路径权重；
各路径内的学习者信念按第\ref{sec:feedback}节的规则更新。

\subsection{以选择为观测的状态推断}
参数向量记为$\boldsymbol\psi$。行为路径使用式\eqref{eq:path-state}中已包含本题执行身份的$Z_t$。
给定该状态，$p_{\boldsymbol\psi}(a_t=c\mid Z_t)=\sum_gq_{tH_t}(g)p_{tH_tg}(c\mid\bm u_t)$。
式\eqref{eq:choice}则是在抽取执行身份之前对它求和的预测。逐题过程写成状态转移核
$\mathcal K_{\boldsymbol\psi}(Z_{t+1}\mid Z_t,a_t,r_t,\bm x_{t+1})$。
这个核先完成本题反馈学习和经历保存，再产生下一题知觉及内容形成，最后依据本题执行身份
抽取下一题执行内容。它在转移中保留执行维持所需的历史。
令$\mu_{\boldsymbol\psi}(Z_1\mid\bm x_1)$为第一题初始状态分布，固定任务日程下的
整条生成路径满足
\begin{equation}
\begin{split}
p_{\boldsymbol\psi}(Z_{1:T},a_{1:T}\mid\bm x_{1:T},\mathcal E)
={}&\mu_{\boldsymbol\psi}(Z_1\mid\bm x_1)
\prod_{t=1}^{T}p_{\boldsymbol\psi}(a_t\mid Z_t)\\
&\times\prod_{t=1}^{T-1}
\mathcal K_{\boldsymbol\psi}(Z_{t+1}\mid Z_t,a_t,r_t,\bm x_{t+1}).
\end{split}
\label{eq:path-model}
\end{equation}
$\mathcal E$表示已知实验环境，按真实类别和当前选择确定$r_t$。
学习者的转移核接收产生的反馈。行为推断对式\eqref{eq:path-model}中的潜在路径求和，
自主模拟则从同一式逐步抽样。
该式把刺激日程视为给定输入；实验按既往反应调整出题及结束任务的概率属于环境过程，
本版对这些概率的建模范围见下文自主生成说明。

知觉、形成和选择共同定义一个具有历史依赖的生成模型。研究者用$N$个粒子近似其潜在状态
分布，每个粒子保留一条自身一致的学习历史。$N$控制计算精度，学习者容量$K$规定每条
路径内能够保留的内容数，两者分别设定与解释。在第$t$题回答之前，每个粒子已抽取本题执行内容，第$i$个粒子给出
相应的$p_t^{(i)}(a_t)$。当$\kappa=0$时，可解析积分执行身份，采用式\eqref{eq:choice}的混合概率。预测当前实际选择并更新路径权重的公式为
\begin{equation}
\widehat P(a_t\mid\mathcal D_{<t},\bm x_t,\boldsymbol\psi)
=\sum_iW_{t,-}^{(i)}p_t^{(i)}(a_t),\qquad
W_{t,a}^{(i)}=\frac{W_{t,-}^{(i)}p_t^{(i)}(a_t)}{\sum_jW_{t,-}^{(j)}p_t^{(j)}(a_t)}.
\label{eq:filter}
\end{equation}
$\mathcal D_{<t}$为此前已观察的刺激、选择和反馈。随后每个粒子接收当前反馈，按同一个
学习算子更新自身状态。有效粒子数为$1/\sum_i(W_{t,a}^{(i)})^2$；低于粒子数的一半时，
使用共享滤波器的重采样步骤。被复制粒子的未来知觉和搜索随机数流重新分配。

用$\widehat{\mathcal L}_T$表示行为似然估计，它是逐题预测实际选择概率的乘积：
\begin{equation}
\widehat{\mathcal L}_T(\boldsymbol\psi)=\prod_{t=1}^T
\widehat P(a_t\mid\mathcal D_{<t},\bm x_t,\boldsymbol\psi),\qquad
\mathrm{NLL}_T=-\log\widehat{\mathcal L}_T.
\label{eq:behavior-likelihood}
\end{equation}
NLL为负对数似然；实际选择获得的预测概率越高，NLL越低。
当前反馈在当前选择评分之后进入学习。实验环境决定反馈，反馈同时构成后续状态的已知输入。
研究者已知实验类别和反馈程序，给定实际刺激与选择后，环境反馈已经确定。
学习者内部的反馈似然用于改变内容信念；研究者的行为似然对实际选择评分。
本轮条件3通过完整三值反馈更新内部状态。

每个固定参数点运行$R$个独立粒子系统，以整段似然的均值近似边际似然。
这里的边际化针对知觉、内容形成及执行路径，个体参数$\boldsymbol\psi$在每次估计中固定。
若前$m$题用于估计参数，则顺序预测后段使用
\begin{equation}
\widehat P(a_{m+1:T}\mid\mathcal D_{1:m},\boldsymbol\psi)
=\frac{R^{-1}\sum_{r=1}^R\widehat{\mathcal L}_T^{(r)}}{R^{-1}\sum_{r=1}^R\widehat{\mathcal L}_m^{(r)}}.
\label{eq:future-likelihood}
\end{equation}
分子与分母来自同组完整滤波运行，因而保留前段对潜在路径的条件化作用。
计算采用对数域的log-sum-exp。重复系统的相对权重和不同粒子数间的差异作为数值诊断。

数值标准误采用独立系统之间的一阶误差传播近似。令
$v_r^T=\widehat{\mathcal L}_T^{(r)}/\sum_j\widehat{\mathcal L}_T^{(j)}$及
$v_r^m=\widehat{\mathcal L}_m^{(r)}/\sum_j\widehat{\mathcal L}_m^{(j)}$，对于$R\geq2$，后段对数似然的标准误为
\begin{equation}
\widehat{\mathrm{SE}}=\left[\frac{R}{R-1}\sum_{r=1}^R(v_r^T-v_r^m)^2\right]^{1/2}.
\label{eq:mc-error}
\end{equation}
每题评分再除以后段题数。重复数有限，该近似的置信覆盖率尚待校准；
数值检查同时保留粒子数变化和单个系统似然质量集中的信息。

\subsection{状态时点与可解释范围}
选择前的预测状态条件于此前的选择与反馈；选择后的过滤状态再条件于当前选择，两个时点
都在当前反馈之前。参数在本轮状态检查中固定，参数不确定性另行比较。
真人参数由训练前段估计，训练段内的状态汇总属于选参后的回顾性描述；后段顺序预测使用
已经固定的参数。模拟状态恢复则直接使用已知生成参数。
在$\kappa=0$的解析边界，观察$a_t$后，执行身份的条件分布在粒子内部更新为：
\begin{equation}
P(H_t=h\mid a_t,S_t,\bm u_t)
=\frac{w_{th}p_{th}(a_t\mid\bm u_t)}{\sum_jw_{tj}p_{tj}(a_t\mid\bm u_t)}.
\label{eq:execution-filter}
\end{equation}
它用于研究者对执行内容的推断。$\kappa>0$的抽样模式把实际执行身份保存在各粒子路径中，
并通过式\eqref{eq:filter}中的粒子权重对它条件化。学习者的内容信念按实际反馈更新。

跨粒子汇总采用具有相同心理含义的内容，包括具体样例、局部概括、目录规则、保留维度和
明确的目录规则编号。样例中心及来源保存在各路径内部。均值和逐题概率描述分布，
一条完整路径描述一次连贯过程。口述是对可报告决策内容的外部有噪观测；本版尚未建立
模板内容到完整口述的生成概率。本轮口述不参与选参、结构选择或新增验证。
Task2实验程序在作答后、本题反馈显示前录音，因此口述相关的候选状态对齐选择后、
当前反馈前的时点。程序的反馈表同时规定了条件3的$1,\tfrac12,0$三种事件。

\subsection{同一生成模型中的自主反馈}
自主模拟固定参数与任务刺激日程，先按式\eqref{eq:execution-policy}抽取执行内容$H_t$，
再按该内容的类别预测抽取选择。式\eqref{eq:choice}给出这两步的选择边缘概率。
实验环境根据真实类别、配对及模拟选择产生反馈；下一题依据这条路径自己的执行身份、
选择和反馈继续学习。
所有路径保留其发生情况，用于估计轨迹形态、改善时间、维持表现及其概率。
本轮固定已记录的刺激日程和记录长度；实验的自适应出题及达标停止机制尚未联合模拟。

模型包含学习者的反馈学习和研究者对生成路径的贝叶斯推断两个层次。
给定当前内容集合及保持后的有效旧分布，式\eqref{eq:one-event}执行一次条件贝叶斯更新。
内容保持指数依赖本题反馈，完整认知更新包含反馈调节的幂保持与随后的条件化两个步骤。
当$\delta=1$时，它退化为对原信念的普通贝叶斯更新。候选形成、容量替换、精度学习及
执行选择共同规定状态转移，研究者对该生成模型的潜在路径进行贝叶斯推断。

\section{版本比较、估计与证据范围}
\label{sec:evidence}

\subsection{被试、数据用途和计算方案}
本轮使用18名被试的9792题有效记录。每人按原有顺序取前70\%用于估计参数，
逐人取整后共6846题；其余2946题用于顺序预测。开发被试的后段参与结构选择，
评价被试的后段选择在原七参数结构冻结后进入预测评分。该检验失败后，这6人的后段
又用于修订比较，因此修订版本的评价属于开发分析。名单按原始分工列出。
\begin{center}
\begin{tabularx}{\textwidth}{p{1.3cm}Y Y}
\toprule
条件 & 开发被试 & 结构冻结后的评价被试\\
\midrule
1 & 110、113、118、128 & 103、132\\
2 & 202、205、215、229 & 208、221\\
3 & 303、306、315、317 & 304、331\\
\bottomrule
\end{tabularx}
\end{center}
这些角色是本轮研究的数据用途划分。既有项目中被查看过的案例继续具有开发历史。
口述资料等待最终人工复核，本轮只使用行为记录和模拟真值。

参数采用有限离散候选的似然搜索。共同初始池含48点，其中16点为此前先导在拟合前生成的
全部预设点，其余来自新Sobol序列；旧结果的排序未参与候选生成。各结构从共同点投影到
自身参数边界并去重，再从领先点生成16个局部候选。每个局部候选随机改动1至3个坐标，
改动值从预定离散集合重选，然后按结构约束投影。既有16点的坐标连同新搜索网格共同
构成本轮候选来源。嵌套大模型在确认前接收小模型的领先点。
初筛使用24粒子、2次重复；4个点以96粒子各2次确认，2个点以256粒子各4次复核，
必要时增加到512粒子各4次。两个预选结构以1024粒子各4次独立检查后段预测。
这些粒子数构成工作精度方案，其实际稳定范围随结果报告。

开发后段首轮比较完成后，提前计算两个可能入选结构各自与完整参考版在评价被试训练段
上的二模型搜索，分别保存完成表。每一对结构的搜索规则、候选和种子保持原值。
原结构冻结后采用相应训练结果，再评分评价后段。失败之后，修订比较使用另一份
已完成并冻结的九参数配对训练结果，重新评价同一批人的后段；参数继续只由前段确定。
预计算与开发复核的数值工作进程总数上限为128，CPU消耗由合并计数器监管。

评价被试结果产生前另登记固定参数的可选复核：原定恢复完成后，若剩余预算可以容纳
同配置实测成本的两倍及五分钟余量，则以1024粒子各4次独立运行检查评价后段。
模型结构和训练选定参数保持固定，实际执行状态与原评分分别记录。

结构选择首先比较开发被试等权的后段每题NLL。工作规则优先保留平均差加入数值余量后
不超过0.01、各条件平均差不超过0.03、自主Energy恶化不超过10\%的简洁结构。
数值余量逐被试计算：两结构的完整配置相同时，比较中的数值误差相消；其余被试取两版本
标准误之和。对被试取平均后再乘二，得到最终余量。阈值及余量是本轮版本决定的工作标准。
通过这些标准的版本先按名义自由参数数排序，再按可选控制机制的保留数量排序，最后比较
后段NLL。计数包括额外执行维持、精度学习、累积失败搜索、反馈调节保持及相似度检索。
同参数数的这一优先顺序在开发后段评分产生前登记。
同种子下不同粒子数的运行先在种子内合并，再跨独立种子估计误差。
独立精度复核中，每结构仍运行4次；相同被试的完整配置若在两个结构中相同，则合并其
8次独立运行供两者共用，并保存实际运行数与来源任务。相同配置的预测差因而为零。
跨被试复用种子时，平均评分的数值误差按各被试标准误的平均作保守合并。
独立重复的整段似然质量集中或粒子数之间差异较大时，保留数值限制标记。
作为评分尺度的参照，逐题均匀随机选择的每题NLL为$\log C$，二分类约为0.6931，
四分类约为1.3863；分数降低表示实际选择获得了更高的预测概率。

本轮参数搜索给出固定参数点及已检查候选的似然排序，完整参数后验仍待估计。
八个受比较版本均含经验工作区，因此本轮决定针对共同表示体系内的控制机制。
经验表示相对于纯规则体系的贡献，以及各形成操作的必要性，需通过对应的表示消融检验。
嵌套检查比较完整模型与其简化边界上已经确认的训练参数点。相同粒子数和重复数下，
边界点若在计入数值误差后仍明显较好，则标记有限搜索不足；不同计算精度的比较分别记录。
早期共享边界候选有助于搜索，后续有限候选筛选仍可能遗漏更好的参数点。
有限搜索、参数补偿和数值误差分别评价。参数在一段任务内固定，稳定心理特质的解释需要
跨任务或重复测量的进一步证据。

\subsection{完整路径的分布评价}
每个被试与所选参数组合在实际刺激日程上自主生成64条路径，全部路径参与评价。
特征向量包含32题块正确率、最大相邻块变化及其时点、该变化之后的平均水平、块间变动、
各正确类别下的反应频率；条件3另加入同科正确率和平均反馈得分。
最大变化允许为负，随后的水平用于说明改变之后的实际表现。

设第$m$条路径的特征向量为$\bm f^{(m)}$，真人特征为$\bm f^{\mathrm{obs}}$，
特征维度为$d$、路径数为$M$。本轮采用越低越好的Energy评分方向\cite{gneiting2007}，
对全部路径组成的经验预测分布计算：
\begin{equation}
\mathrm{ES}=\frac1M\sum_m\frac{\|\bm f^{(m)}-\bm f^{\mathrm{obs}}\|_2}{\sqrt d}
-\frac1{2M^2}\sum_{m,n}\frac{\|\bm f^{(m)}-\bm f^{(n)}\|_2}{\sqrt d}.
\label{eq:energy}
\end{equation}
分数越低表示该特征空间中的预测分布越接近观察。第一项评价预测与观察的距离，第二项
计入预测分布自身的离散程度。该检查使用完整日程，包含训练段；后段选择预测单独报告。
所有路径各自保持一条连续学习历史，路径均值和区间描述其分布。
结果另列各特征的5\%、50\%和95\%分位数。联合特征覆盖以模拟特征均值为中心，
取模拟路径到中心距离的90\%分位数作为半径，报告真人特征是否落在该区域及模拟质量。
这些摘要描述固定参数下完整日程的经验预测分布；结构选择使用上述预定Energy评分。
另按种子顺序把64条路径分成两个不重叠的32条路径子集，比较Energy、平均正确率及内容组
占比，描述有限模拟次数下的变化。正式结构选择继续使用全部64条路径。

\subsection{模拟真值与恢复}
有限恢复使用S110、S202、S303的刺激与知觉背景，原初选七参数版与完整十参数版
各生成2条路径，共12条。恢复对象在保留样本反证出现前确定，其范围按这两个版本记录。
这三个背景均采用阈值均匀知觉误差；正态误差背景的恢复范围仍待评价。
每条路径以64个新Sobol点及此前固定的全部16个预设点重新搜索，按结构约束投影并去重，
每种结构的前两点以256粒子各2次确认。
生成参数另用独立种子评分，用于区分搜索不足和机制难以区分。
预设点补充在开发首轮比较完成后、恢复资料生成前登记，用于覆盖既有候选中的容量10等网格外坐标。
原64个新候选和种子保留，补充采用完整预设集合。
在生成结构下，逐坐标报告生成值、恢复值及误差，并标出结构固定项与生成点上不起作用的
坐标。这些检查描述本轮有限生成点，完整可接受参数区域仍需另行估计。
该比较使用完整模拟记录的似然，解释为有限候选中的诊断。边界上配置等价的生成点
按最简有效配置记录。
等价检查纳入容量为1时的执行维持、检索和内容保持方式，以及恒定搜索时不起作用的累积失败保持。
条件1、2容量为1时的保持率也失去作用；条件3的同一参数仍控制配对信念。

状态恢复固定生成参数，以512粒子各2次推断选择之后、当前反馈之前的状态。
目标包括三类内容及初始占位的执行概率、内容组信念，以及各特征上的模板信念质量，
后者等于保留该特征的模板信念之和。条件3另评价配对边缘信念。
逐题执行概率用实际采用的内容组真值评价，目录规则身份单独报告其真值获得的概率。
内容组信念报告均值误差和边缘90\%区间覆盖；模板维度及配对信念报告均值误差。
配对状态的恢复目标是模拟学习者当时的信念；环境的真实配对用于生成反馈和评价任务表现。
覆盖率汇总来自相关试次，其作用是检查已知模型和参数下的近似恢复。
模拟提供完整潜在真值，真人口述在后续工作中提供有噪外部观测。

状态恢复输出产生前，另登记一个描述性参照：使用生成参数下已经保存的64条自主路径，
以各时点的内容组频率表示尚未观察行为时的状态预测。对内容组采用Brier评分，即预测概率
与实际执行组的单点指示向量之间的平方误差之和；评分越低越好。
将它与观察选择后的过滤评分比较，描述行为提供的状态信息。
这64条路径也用于开发阶段的自主评价；恢复真值路径采用另一组种子。
该追加汇总使用既有数组，模型选择和恢复分类保持原定义。
这些目标检验具体的状态摘要。记忆对象身份、模板中心及完整联合状态的恢复仍待评价。

\subsection{本轮结果}
\subsubsection{结构选择与开发证据}
开发阶段首先选中七参数候选，随后保留样本的评价推翻了该采用决定。
以下先记录原有选择依据，再报告反证与修订比较。原有结构冻结文件和评价结果完整保留。

八个版本的开发后段评分如下。各人等权，NLL和Energy均越低越好；
NLL来自首轮固定参数评价，自主Energy使用各组全部64条路径。
\begin{center}
\begin{tabularx}{\textwidth}{Y r r r}
\toprule
版本 & 参数数 & 后段每题NLL & 自主Energy\\
\midrule
完整参考 & 10 & 0.522459 & 0.124722 \\
取消额外执行维持 & 9 & 0.521740 & 0.130892 \\
固定共享精度 & 9 & 0.533736 & 0.131097 \\
共同简化 & 8 & 0.526773 & 0.137700 \\
共同简化及恒定搜索 & 6 & 0.541559 & 0.146694 \\
共同简化及恒定保持 & 8 & 0.527013 & 0.147941 \\
七参数初选候选 & 7 & 0.525046 & 0.125689 \\
共同简化及相似度检索 & 8 & 0.531001 & 0.145469 \\
\bottomrule
\end{tabularx}
\end{center}

共同简化后进一步改为恒定搜索时，预测差与自主分布差超出本轮容忍，因此保留累积失败搜索。
相似度检索和恒定保持版本也未通过全部工作标准。各形成操作及每个剩余自由坐标的独立必要性
继续按对应消融与恢复问题评价。

两个预选结构随后采用1024粒子、各4个独立系统复核；3名被试的配置相同，分别合并为8个系统
供两者共用。下表的数值标准误描述计算波动，被试间差异另行保留。
\begin{center}
\begin{tabularx}{\textwidth}{Y r r r}
\toprule
版本 & 后段每题NLL & 数值标准误 & 数值受限条目\\
\midrule
取消额外执行维持 & 0.522043 & 0.000689 & 2/12 \\
七参数初选候选 & 0.526035 & 0.000550 & 2/12 \\
\bottomrule
\end{tabularx}
\end{center}

七参数版的平均NLL高0.003993；计入数值余量后的差为0.005584，低于0.01工作门槛。
分条件最大差为0.010817，低于0.03；自主Energy降低3.98\%。
因此，原定规则在这一阶段选中七参数候选。三个条件使用共同参数定义。

七参数版中，S215的独立系统似然质量高度集中，S317的不同粒子数评分差超出本轮数值标准；
九参数版的受限条目为S113和S317。S317在两版本中使用相同配置，结构差值使用同一合并估计，
绝对数值精度仍保留限制。个体参数可接受区域、稳定心理特质和完整内部状态准确性分别需要证据。

18名被试的七参数训练估计中，7例保留训练评分数值限制；11例的第二候选与领先候选的
原始训练NLL差不超过每题0.01。这里的差值描述已检查候选的接近程度，参数区间的覆盖率
需另行估计。S306、S317、S304取$\eta=0$，对应的$\rho$不起作用；5例容量取本轮上限10，
1例判断精度取上限25。这些边界与候选间差异随个体结果保留。

\subsubsection{初选方案的保留样本检验}
原先冻结的七参数版在6名评价被试上未达到预定的相对预测容忍范围。参数均由前70\%记录确定，
下表评分来自此前保留的后30\%；各人等权。
\begin{center}
\begin{tabularx}{\textwidth}{Y r r r r r}
\toprule
版本 & 总体NLL & 条件1 & 条件2 & 条件3 & Energy\\
\midrule
七参数初选 & 0.557058 & 0.339383 & 0.551686 & 0.780104 & 0.147419\\
完整十参数 & 0.376043 & 0.296869 & 0.437470 & 0.393791 & 0.127031\\
\bottomrule
\end{tabularx}
\end{center}
七参数版的每题NLL高0.181015，自主Energy高16.05\%。条件3的平均差为0.386314。
S103、S208、S221、S304、S331的NLL分别高0.0850、0.2032、0.0253、0.6035和0.1691；
S132在两结构中采用同一配置，差为零。这些结果超过原有精简容忍范围，
因此撤回七参数版的采用决定。它同时固定了三个控制坐标，失败本身尚不能把损失归于其中某一个。

原评价中S304的两版本训练估计均保留数值限制，完整版本的后段独立系统质量也较集中。
可选1024粒子复核未启动：按登记的两倍实测成本规则估计需4.94小时，
当时剩余计算时间为1.59小时。原评价的256或512粒子、4次独立结果完整保留。
各人原始差值与数值诊断共同限定解释范围；明显的多被试差距足以否定本轮七参数采用方案，
精确的总体效应大小仍需更多独立数据和计算精度检验。

\subsubsection{修订比较与本轮采用决定}
原检验失败后，采用事先已经完成的九参数版本与完整版本的配对训练表，评价同一批6人的后段。
参数保持由前70\%选择确定，评分方法、随机种子、粒子数及每配置64条自主路径沿用原方案。
完全相同的已完成任务直接复用。由于修订决定已经受到这6人的原评价影响，本次结果属于
开发分析；修订后的结构尚需在另一批独立资料上验证。
\begin{center}
\begin{tabularx}{\textwidth}{Y r r r r r}
\toprule
版本 & 总体NLL & 条件1 & 条件2 & 条件3 & Energy\\
\midrule
取消额外执行维持 & 0.404714 & 0.330446 & 0.438237 & 0.445459 & 0.142717 \\
完整十参数 & 0.379930 & 0.302784 & 0.438307 & 0.398699 & 0.126329 \\
\bottomrule
\end{tabularx}
\end{center}
九参数版与本次完整版的平均NLL差为0.024784。相对本次最低NLL版本，加入数值余量后的差为0.027191，
分条件最大差为0.046760，自主Energy相对差为12.97\%。
九参数方案未通过预定的全部保护标准，因此本轮采用完整十参数版本。
具体样例、局部概括与目录规则继续共享有限工作区、反馈事件和内容信念；本轮保留
精度学习、信念保持和执行维持三个控制。它们分别调节预测强度、旧证据优势和执行连续性。
本次结果支持在已比较版本中保留该结构；每个坐标的独立必要性和最佳有效维度继续按
更充分的参数搜索、独立行为资料与恢复证据判断。
追加比较中，S304的两版本训练评分均受数值精度限制。九参数版的S221、S304、S331，
以及完整版的S304，其后段独立系统似然质量集中在一个系统上的比例超过0.8。
这些标记与工作门槛一同保留，独立重复的近似标准误对数值不确定性的刻画仍有限。

完整版本18名被试的领先训练点中，15人的$\kappa=0$；S118与S303取6，S304取3。
因此，该参数在当前点估计中的作用集中于少数案例。完整模型有7例训练评分数值受限，
14例的第二候选与领先候选相差不超过每题0.01；5人的容量达到本轮上限10。
这些结果说明，本轮采用决定与可靠的个体参数识别仍有明显距离。

\paragraph{既有零维持边界候选的敏感性。}
修订比较中，除S304外，完整版本的5个领先点都落在$\kappa=0$边界，因而也是合法的九参数配置。
其中部分点与九参数原搜索的领先点不同，表现出有限搜索的影响。使用已经计算的这些边界点，
按训练NLL转入九参数比较，九参数的后段平均NLL变为0.397058，与完整版本相差0.017128；
计入数值余量后为0.017656，条件3的差为0.051383，仍超过相应工作门槛。
S103的两份训练评分使用不同粒子数，此转入结果按敏感性描述。
该检查使用既有数组，登记比较与最终配置保持原值。

另作一个对已知零维持候选有利的描述性检查：每人从本次修订比较的两份领先配置中，
直接取属于零维持边界且未来段NLL较低的一份，
所得均值为0.395493，仍比完整版高0.015563。该最小值使用了未来选择，只界定这组有限候选的
乐观评分，个体选参与预测评价继续使用原登记结果。
这些检查表明，纳入上述已计算边界点后，本轮采用决定保持一致；剩余差距主要
集中于S304，该例本身仍有数值限制。执行维持的独立必要性继续需要更充分搜索与独立资料。

两次评价中的完整版本分别来自各自的配对训练搜索，部分被试的参数点不同，因此两表的
完整版分数分别报告。参数表及配置身份随结果保存，未来段分数未用于重新挑选个体参数。
原冻结决定保存在\path{structure_selection.json}，修订方案与结果分别保存在
\path{revision_plan.json}、\path{revision_selection.json}及\path{report_v2/analysis.json}。


\subsubsection{参数与结构恢复的限制}
12条模拟路径的结构恢复分类为：1条能明确识别，3条尚难区分，8条受数值精度限制。
11条的简洁工作选择与生成配置的有效结构一致；该比例包含数值受限案例，
其作用是描述有限搜索结果，结构识别率的可靠估计仍待校准。
S110的4条路径虽然分别挂在两个名义结构下，实际生成配置完全相同，均落在七参数边界上。
这部分只能检查共同边界的恢复。唯一明确区分的案例来自S202的完整模型生成路径。

逐坐标恢复保留了明显误差。完整有效参数向量仅3/12条精确恢复，均来自S110背景；
容量恢复平均绝对误差为1.58，相似宽度为0.113，规则提议概率为0.063。
七参数及完整版本共有坐标的精确恢复分别为：容量7/12、相似宽度6/12、初始精度7/12、
基准搜索3/12、失败增益6/12、有效失败保持3/12、规则提议10/12。
完整版本另外三个有效坐标各仅4条生成路径，保持率、执行维持和精度学习率均为2/4精确恢复。
生成参数评分和重复系统诊断随每条路径保存，有限搜索不足、参数补偿与数值波动共同影响这些结果。
本轮据此保留点估计与候选差异，个体参数的完整可接受区域和跨任务稳定性继续待检验。

\subsubsection{固定生成参数下的状态恢复}
下表汇总已知生成参数时，选择后、当前反馈前的状态恢复。
“真组概率”是后验分配给该题实际执行内容组的平均概率；“真规则概率”在实际执行目录规则的
题目上，对具体规则编号作相同计算。二者都描述真值得到的概率质量。
\begin{center}
\begin{tabularx}{\textwidth}{Y r r r r r}
\toprule
生成版本与条件 & 真组概率 & 真规则概率 & 组信念MAE & 90\%覆盖 & Brier\\
\midrule
七参数，条件1 & 0.806 & 0.774 & 0.056 & 0.866 & 0.196\\
完整版，条件1 & 0.843 & 0.849 & 0.049 & 0.914 & 0.155\\
七参数，条件2 & 0.696 & 0.192 & 0.073 & 0.953 & 0.290\\
完整版，条件2 & 0.781 & 0.447 & 0.038 & 0.959 & 0.223\\
七参数，条件3 & 0.830 & 0.617 & 0.053 & 0.856 & 0.178\\
完整版，条件3 & 0.894 & 0.597 & 0.036 & 0.959 & 0.100\\
\bottomrule
\end{tabularx}
\end{center}
每行对应2条路径、各2个数值系统。两版本使用不同生成路径与参数，表中差异描述各自的
恢复范围。MAE为平均绝对误差；90\%覆盖为各内容组信念的边缘区间覆盖率，
其中包含真实质量为零的组。分组及正质量试次的覆盖单独保存在恢复审查结果中。
条件3的学习者配对信念平均绝对误差为0.029（七参数）和0.040（完整版）。
同一路径两个数值系统的内容组执行概率平均差范围为0.0028--0.0152；
这些摘要的重复性与完整联合状态的准确恢复分别评价。

与64条自主路径构成的无行为条件化参照相比，过滤后的内容组Brier在六个“生成版本\(\times\)条件”
组合中均降低。七参数的降幅依次为0.0332、0.0574、0.0190，完整版为0.0032、0.0209、0.0174。
这说明在这些模拟例子中，观察选择为内容组推断提供了额外信息；该比较使用有限生成点及已知参数。
条件2的具体规则身份仍较难恢复。本轮未专门评价“正确规则在场但未执行”这一状态，
既有执行维持解释的恢复边界继续保留。记忆对象身份、中心位置、参数积分后的状态分布
以及真人口述的生成概率继续留待专门检验。


\section{实现、复现与使用边界}
\label{sec:implementation}
模型由共享的\texttt{StateModel}实现，三条件使用共同的内容形成、选择和反馈时序。
条件1、2原有v1配置继续保留；统一版本为\texttt{1031-unified-v2}。
三份最终配置位于\path{configs/exp123/model_struct/}，文件名分别为
\path{model_1031_cond1_final_v2.yaml}、
\path{model_1031_cond2_final_v2.yaml}和
\path{model_1031_cond3_final_v2.yaml}。
这些配置采用十参数完整结构，数值初值作为接口示例使用；
各被试的训练估计及复核候选随研究结果保存。
示例配置的执行维持为正，生成与推断均保留抽样的执行身份。
仅在$\kappa=0$边界，行为推断可把\path{marginalize_execution}设为\texttt{true}，
按式\eqref{eq:choice}逐题精确求和执行身份。
配置中的16粒子用于小规模接口检查，正式推断按对应研究的数值精度方案设置。
载入最终YAML并传入被试编号与条件后，三条件可通过共享\texttt{StateModel}运行；
本轮的\path{model_1031_finalization}研究入口负责数据准备与拟合。
早期\path{run_model_1031_smoke}入口的适用范围仍为条件1、2。

本轮实现检查覆盖联合反馈恒等式、精度梯度、概率归一化、因果时序、状态快照及旧版数值回归。
最终三份配置各用24题有效合成序列运行生成与重放，预测概率最大差均为0；
再各用8粒子检查选择前与选择后的状态。条件3覆盖全错、半对、完全正确三类事件，
内容信念、配对分布及条件类别分布均保持归一化，判断精度保持在规定区间。
全部17264份数值任务记录的数组哈希通过核验，16620份拟合记录从保存概率重算NLL的最大差为0，
完整数组均有对应任务记录。上述检查说明实现与记录的一致性，模型预测与状态恢复结果见第\ref{sec:evidence}节。

研究配置、输入数组、知觉汇总及源码哈希，随机种子、逐任务输出和中断记录随研究目录保存。
数值核心冻结于提交\texttt{6c3a4be9}，运行时逐文件核对登记哈希；协调与审查代码另保存阶段版本。
本轮目录为\path{results/model_1031/finalization_20261101_v2/}。旧结果和原始数据保持原状。
在仓库根目录用\texttt{bash}运行
\path{src/Bayesian_state/docs/model_architecture/compile_model_1031.sh}，
并把新PDF路径作为参数，即可重新编译本文。脚本在临时目录完成三次编译，已有目标文件保持原状。

工作区的数学状态按具体内容定义。实现中的存储位置用于访问数组，位置替换不会延续旧内容
的执行身份。模板中心、保留维度、类别信念、条件配对及来源关系随粒子状态一起保存和恢复。
历史预测使用对应试次的内容与配对重放。
粒子快照在相同配置下恢复；跨参数的状态推断分别运行相应配置。
内部数组和试次日志从0开始索引，本文的第$t$题从1开始编号。

抽样生成中的\texttt{choice\_weights}记录本题已抽到内容的单点权重，抽样之前的内容选择
分布保存在转移日志的\texttt{execution\_selection\_probabilities}中。
解析推断模式的\texttt{choice\_weights}保存完整执行概率；代表编号取最大执行概率项，
其持续计数描述该代表编号的连续出现。抽样执行身份的连续路径由自主生成记录。
固定精度配置把\texttt{learning\_rate}设为0，模板类别信念仍在反馈后按
式\eqref{eq:label-update}更新。

目录的完整区域系数由\path{src/Bayesian_state/hypothesis_space/spaces/continuous.py}和
同目录的\path{structural_0923.py}生成。四分类扩展标识为\texttt{relational\_152\_0923}；
维度间相近带和单维中心带的固定容差均为0.10。这些目录常数与模板的相似宽度$\sigma$
分别定义：前者规定可供选择的规则内容，后者控制经验模板的适用范围。
研究目录中的\texttt{catalogue\_audit.json}保存目录内容签名、规则族数量及重复检查，
\texttt{task\_catalogue\_audit.json}保存记录刺激的逐规则分类覆盖。

选择前与选择后、反馈前的内容状态分别标记。运行输出另保留反馈后的配对边缘：
\texttt{pairing\_prior}为选择前状态，\texttt{pairing\_posterior}为当前反馈之后的状态。
反馈前的选择后配对状态由粒子观察接口给出。本轮状态恢复使用该接口对齐模拟真值。

模型假定有限工作区、固定模板中心、逐次删除一个维度、固定局部操作比例及预设规则目录。
这些内容构成明确的认知假设。目录几何重复检查和数值检验支持实现的一致性；
表示资源的完整性、各形成操作的必要性、完整参数区域及口述生成机制仍有各自的验证任务。
固定刺激日程下的路径预测也保留对自适应出题和任务停止机制的条件化范围。

\appendix
\section{受比较的控制机制}
\label{sec:variants}
完整参考版有十个自由坐标：$K,\sigma,\beta_0,p_0,\eta,\rho,g_{\mathrm R},\delta,\kappa,\alpha$。
以下说明与最终版本比较的控制机制。
完整版本按式\eqref{eq:execution-policy}保留执行维持。取消额外维持的比较版本固定$\kappa=0$，
行为推断可逐题解析求和执行身份，其余认知状态与学习步骤沿用共同模型。

精度学习见式\eqref{eq:selected-precision}。固定共享精度版本取$\alpha=0$，
共同初值与新内容的精度继承使各项内容保持同一$\beta$。模板类别信念继续接受反馈更新。

累积搜索使用式\eqref{eq:search-pressure}和式\eqref{eq:search-probability}。
恒定搜索版本取$\eta=0$，此时$\rho$对行为预测失去作用。

保持比较包括$\delta_t=\delta+(1-\delta)y_t$的反馈调节形式、$\delta_t=\delta$的恒定形式，
以及$\delta=1$的完全保持。相似度检索版本使用$w_h\propto\pi_h A_h(\bm u_t)$，
规则的检索亲和度固定为1。相似度同时影响模板类别预测与当前内容检索，属于独立的选择假设。

\begin{center}
\begin{tabularx}{\textwidth}{Y p{2.2cm} p{2.3cm} L{1.6cm}}
\toprule
版本 & 执行维持 & 精度 & 参数数\\
\midrule
完整参考 & 自由$\kappa$ & 学习$\alpha$ & 10\\
取消额外维持 & $\kappa=0$ & 学习$\alpha$ & 9\\
固定共享精度 & 自由$\kappa$ & 固定$\beta$ & 9\\
共同简化 & $\kappa=0$ & 固定$\beta$ & 8\\
共同简化及恒定搜索 & $\kappa=0$ & 固定$\beta$ & 6\\
共同简化及恒定保持 & $\kappa=0$ & 固定$\beta$ & 8\\
共同简化及完全保持 & $\kappa=0$ & 固定$\beta$ & 7\\
共同简化及相似度检索 & $\kappa=0$ & 固定$\beta$ & 8\\
\bottomrule
\end{tabularx}
\end{center}
名义维度描述结构的自由坐标数。特定边界的有效维度可以更低，例如$K=1$时额外执行维持
不影响选择，$\eta=0$时失败保持系数失去作用。条件1、2在$K=1$时内容信念恒为1，
内容保持率也不影响行为；条件3的保持率仍影响配对信念。这些边界按有效参数解释。
高精度下模板增益趋于饱和，$\sigma$、类别
不确定性与内容信念也可能形成实际估计中的补偿。结构简洁性与个体参数可辨识性分别评价。

\begin{thebibliography}{9}
\bibitem{nosofsky1994}
Nosofsky, R. M., Palmeri, T. J., \& McKinley, S. C. (1994).
Rule-plus-exception model of classification learning.
\textit{Psychological Review}, 101(1), 53--79.
\href{https://doi.org/10.1037/0033-295X.101.1.53}{doi:10.1037/0033-295X.101.1.53}.

\bibitem{sanborn2010}
Sanborn, A. N., Griffiths, T. L., \& Navarro, D. J. (2010).
Rational approximations to rational models: Alternative algorithms for category learning.
\textit{Psychological Review}, 117(4), 1144--1167.
\href{https://doi.org/10.1037/a0020511}{doi:10.1037/a0020511}.

\bibitem{gneiting2007}
Gneiting, T., \& Raftery, A. E. (2007).
Strictly proper scoring rules, prediction, and estimation.
\textit{Journal of the American Statistical Association}, 102(477), 359--378.
\href{https://doi.org/10.1198/016214506000001437}{doi:10.1198/016214506000001437}.
\end{thebibliography}

\end{document}
